% preamble.tex
% -----------------------------------------------------------------------
% "Aprendo los números" -- cuaderno diario de números (español, A4).
% Todo el aparato tipográfico vive aquí: paquetes, colores, cajas y las
% macros que usa cada página. Es el motor de "Aprendo a leer"
% (https://github.com/konradcinkusz/learning-to-read) -- la misma letra,
% la misma paleta, la misma cabecera de cada día, las mismas cajas --,
% con lo que cambia al contar en vez de leer: la caja de arriba es "El
% número de hoy" (el número, su nombre, el marco de diez y las cosas que
% se cuentan), y las actividades son las de contar (ver notes/01-plan.md).
%
% Requiere LuaLaTeX (fontspec): la letra es Andika (SIL, OFL,
% fonts/andika/), pensada para primeros lectores -- y sus números son los
% mismos que se trazan en la actividad "Traza" (content/numeros-trazo.json
% sale del mismo glifo, ver tools/gen_numeros_puntos.py).
% -----------------------------------------------------------------------

\usepackage{fontspec}
\setmainfont{Andika}[
  Path = fonts/andika/,
  Extension = .ttf,
  UprightFont = *-Regular,
  BoldFont = *-Bold,
  ItalicFont = *-Italic,
  BoldItalicFont = *-BoldItalic,
]
\setsansfont{Andika}[
  Path = fonts/andika/,
  Extension = .ttf,
  UprightFont = *-Regular,
  BoldFont = *-Bold,
  ItalicFont = *-Italic,
  BoldItalicFont = *-BoldItalic,
]

\usepackage[spanish]{babel}
% babel-spanish hace de " un carácter activo; las comillas escritas a mano
% en el texto se estropean sin esto (ver el mismo comentario, más largo,
% en el preamble.tex de "Aprendo a leer"). Solo funciona así, al empezar
% el documento.
\AtBeginDocument{\shorthandoff{"}}

% El pie de página ("Día N / 260") vive en el pie de verdad, no en la
% página de texto: la caja de actividad llena todo lo que queda de página
% (ver `cajadia`), y un pie dentro del texto la empujaría a otra página.
\usepackage[a4paper,top=18mm,bottom=22mm,left=22mm,right=18mm,footskip=10mm]{geometry}

\usepackage{xcolor}
\usepackage[most]{tcolorbox}
\usepackage{tabularx}
\usepackage{booktabs}
\usepackage{enumitem}
\usepackage{tikz}
\usetikzlibrary{arrows.meta}
\usepackage{graphicx}
\usepackage{multicol}

\pagestyle{empty}
\setlength{\parindent}{0pt}
\setlength{\parskip}{2pt}

% --- cadenas de texto ----------------------------------------------------
% lang/es.tex
% -----------------------------------------------------------------------
% Todas las cadenas de texto del cuaderno, en un solo sitio, como en
% "Aprendo a leer": cambiar una palabra aquí la cambia en las 260 páginas
% a la vez, sin tocar el contenido generado.
% -----------------------------------------------------------------------

% Total de días del cuaderno. tools/gen_numeros.py --check comprueba que
% coincide con su propio TOTAL_DIAS.
\newcommand{\totaldias}{260}

% La cabecera de cada día (\cabeceraDia, preamble.tex).
\newcommand{\lblFecha}{Fecha}
\newcommand{\lblNotas}{Notas}
\newcommand{\lblDia}{Día}
\newcommand{\lblSemana}{Semana}
\newcommand{\lblTrimestre}{Trimestre}
\newcommand{\lblHecho}{Hecho}
\newcommand{\lblHechoSola}{sola}
\newcommand{\lblHechoConAyuda}{con ayuda}
\newcommand{\lblHechoConDificultad}{con dificultad}

% La caja de arriba, la misma todos los días.
\newcommand{\lblNumeroDeHoy}{El número de hoy}

% Los títulos de las cajas de actividad.
\newcommand{\lblTraza}{Traza}
\newcommand{\lblColorea}{Colorea}
\newcommand{\lblRodea}{Rodea}
\newcommand{\lblCuenta}{Cuenta}
\newcommand{\lblBusca}{Busca}
\newcommand{\lblUne}{Une}
\newcommand{\lblCompleta}{Completa}
\newcommand{\lblDecena}{Diez y más}
\newcommand{\lblCompara}{Compara}
\newcommand{\lblVecinos}{Antes y después}
\newcommand{\lblRecta}{La recta}
\newcommand{\lblOrdena}{Ordena}
\newcommand{\lblJunta}{Junta}
\newcommand{\lblSuma}{Suma}
\newcommand{\lblResta}{Resta}
\newcommand{\lblParte}{Parte en dos}
\newcommand{\lblDiez}{Hasta el 10}
\newcommand{\lblProblema}{Problema}
\newcommand{\lblBloques}{Decenas y unidades}
% "Number names" solo sale en "First Numbers" (lang/en.tex), pero su caja
% es la misma en los dos cuadernos (preamble.tex).
\newcommand{\lblNombres}{Los nombres de los números}
\newcommand{\lblTabla}{La tabla del 100}
\newcommand{\lblDinero}{El dinero}
\newcommand{\lblHora}{La hora}
\newcommand{\lblMide}{Mide}
\newcommand{\lblDibuja}{Dibuja}
\newcommand{\lblRepasa}{Repasa}

% Las instrucciones: las lee un adulto en voz alta (van pequeñas y en
% gris, como en "Aprendo a leer"). Las de "colorea", "rodea", "cuenta",
% "une" y "completa" las compone tools/gen_numeros.py con el número y las
% cosas de cada día ("Colorea 2 pelotas.").
\newcommand{\lblInstruccionTraza}{Repasa los puntos con el dedo, y luego con un lápiz. Después, escríbelo tú en la línea.}
\newcommand{\lblInstruccionBusca}{Rodea cada uno que encuentres. ¿Cuántos hay?}

% La tabla de las decenas (D) y las unidades (U) de la caja de arriba.
\newcommand{\lblDecenasCorto}{D}
\newcommand{\lblUnidadesCorto}{U}

% La portada.
\newcommand{\lblTituloPortada}{Aprendo los números}
\newcommand{\lblSubtituloPortada}{Un número nuevo cada semana}
\newcommand{\lblSubtituloPortadaMetodo}{Cuaderno diario para contar, trazar los números y empezar a sumar -- con la familia de \emph{Aprendo a leer}}
\newcommand{\lblPortadaNombre}{Nombre}
\newcommand{\lblPortadaCurso}{Curso}

% La medalla del final de cada trimestre (T1-T3) y el cartel de cada
% diez páginas.
\newcommand{\lblMedalla}[1]{¡Medalla de #1!}
\newcommand{\lblMedallaAnimo}{¡Sigue así, \rule{55mm}{0.4pt}!}
\newcommand{\lblPaginasHechas}[1]{¡#1 páginas hechas!}

% La clave de respuestas.
\newcommand{\lblClave}{Clave de respuestas}


% =========================================================================
% Paleta
% =========================================================================
% La de "Aprendo a leer": cada actividad tiene un color saturado (borde y
% título) y uno claro (fondo), y el color es un canal EXTRA, nunca el
% único -- el título en negrita ya dice qué actividad es. main-bw.tex fija
% \bookcolor{bw} antes de % preamble.tex
% -----------------------------------------------------------------------
% "Aprendo los números" -- cuaderno diario de números (español, A4).
% Todo el aparato tipográfico vive aquí: paquetes, colores, cajas y las
% macros que usa cada página. Es el motor de "Aprendo a leer"
% (https://github.com/konradcinkusz/learning-to-read) -- la misma letra,
% la misma paleta, la misma cabecera de cada día, las mismas cajas --,
% con lo que cambia al contar en vez de leer: la caja de arriba es "El
% número de hoy" (el número, su nombre, el marco de diez y las cosas que
% se cuentan), y las actividades son las de contar (ver notes/01-plan.md).
%
% Requiere LuaLaTeX (fontspec): la letra es Andika (SIL, OFL,
% fonts/andika/), pensada para primeros lectores -- y sus números son los
% mismos que se trazan en la actividad "Traza" (content/numeros-trazo.json
% sale del mismo glifo, ver tools/gen_numeros_puntos.py).
% -----------------------------------------------------------------------

\usepackage{fontspec}
\setmainfont{Andika}[
  Path = fonts/andika/,
  Extension = .ttf,
  UprightFont = *-Regular,
  BoldFont = *-Bold,
  ItalicFont = *-Italic,
  BoldItalicFont = *-BoldItalic,
]
\setsansfont{Andika}[
  Path = fonts/andika/,
  Extension = .ttf,
  UprightFont = *-Regular,
  BoldFont = *-Bold,
  ItalicFont = *-Italic,
  BoldItalicFont = *-BoldItalic,
]

\usepackage[spanish]{babel}
% babel-spanish hace de " un carácter activo; las comillas escritas a mano
% en el texto se estropean sin esto (ver el mismo comentario, más largo,
% en el preamble.tex de "Aprendo a leer"). Solo funciona así, al empezar
% el documento.
\AtBeginDocument{\shorthandoff{"}}

% El pie de página ("Día N / 260") vive en el pie de verdad, no en la
% página de texto: la caja de actividad llena todo lo que queda de página
% (ver `cajadia`), y un pie dentro del texto la empujaría a otra página.
\usepackage[a4paper,top=18mm,bottom=22mm,left=22mm,right=18mm,footskip=10mm]{geometry}

\usepackage{xcolor}
\usepackage[most]{tcolorbox}
\usepackage{tabularx}
\usepackage{booktabs}
\usepackage{enumitem}
\usepackage{tikz}
\usetikzlibrary{arrows.meta}
\usepackage{graphicx}
\usepackage{multicol}

\pagestyle{empty}
\setlength{\parindent}{0pt}
\setlength{\parskip}{2pt}

% --- cadenas de texto ----------------------------------------------------
% lang/es.tex
% -----------------------------------------------------------------------
% Todas las cadenas de texto del cuaderno, en un solo sitio, como en
% "Aprendo a leer": cambiar una palabra aquí la cambia en las 260 páginas
% a la vez, sin tocar el contenido generado.
% -----------------------------------------------------------------------

% Total de días del cuaderno. tools/gen_numeros.py --check comprueba que
% coincide con su propio TOTAL_DIAS.
\newcommand{\totaldias}{260}

% La cabecera de cada día (\cabeceraDia, preamble.tex).
\newcommand{\lblFecha}{Fecha}
\newcommand{\lblNotas}{Notas}
\newcommand{\lblDia}{Día}
\newcommand{\lblSemana}{Semana}
\newcommand{\lblTrimestre}{Trimestre}
\newcommand{\lblHecho}{Hecho}
\newcommand{\lblHechoSola}{sola}
\newcommand{\lblHechoConAyuda}{con ayuda}
\newcommand{\lblHechoConDificultad}{con dificultad}

% La caja de arriba, la misma todos los días.
\newcommand{\lblNumeroDeHoy}{El número de hoy}

% Los títulos de las cajas de actividad.
\newcommand{\lblTraza}{Traza}
\newcommand{\lblColorea}{Colorea}
\newcommand{\lblRodea}{Rodea}
\newcommand{\lblCuenta}{Cuenta}
\newcommand{\lblBusca}{Busca}
\newcommand{\lblUne}{Une}
\newcommand{\lblCompleta}{Completa}
\newcommand{\lblDecena}{Diez y más}
\newcommand{\lblCompara}{Compara}
\newcommand{\lblVecinos}{Antes y después}
\newcommand{\lblRecta}{La recta}
\newcommand{\lblOrdena}{Ordena}
\newcommand{\lblJunta}{Junta}
\newcommand{\lblSuma}{Suma}
\newcommand{\lblResta}{Resta}
\newcommand{\lblParte}{Parte en dos}
\newcommand{\lblDiez}{Hasta el 10}
\newcommand{\lblProblema}{Problema}
\newcommand{\lblBloques}{Decenas y unidades}
% "Number names" solo sale en "First Numbers" (lang/en.tex), pero su caja
% es la misma en los dos cuadernos (preamble.tex).
\newcommand{\lblNombres}{Los nombres de los números}
\newcommand{\lblTabla}{La tabla del 100}
\newcommand{\lblDinero}{El dinero}
\newcommand{\lblHora}{La hora}
\newcommand{\lblMide}{Mide}
\newcommand{\lblDibuja}{Dibuja}
\newcommand{\lblRepasa}{Repasa}

% Las instrucciones: las lee un adulto en voz alta (van pequeñas y en
% gris, como en "Aprendo a leer"). Las de "colorea", "rodea", "cuenta",
% "une" y "completa" las compone tools/gen_numeros.py con el número y las
% cosas de cada día ("Colorea 2 pelotas.").
\newcommand{\lblInstruccionTraza}{Repasa los puntos con el dedo, y luego con un lápiz. Después, escríbelo tú en la línea.}
\newcommand{\lblInstruccionBusca}{Rodea cada uno que encuentres. ¿Cuántos hay?}

% La tabla de las decenas (D) y las unidades (U) de la caja de arriba.
\newcommand{\lblDecenasCorto}{D}
\newcommand{\lblUnidadesCorto}{U}

% La portada.
\newcommand{\lblTituloPortada}{Aprendo los números}
\newcommand{\lblSubtituloPortada}{Un número nuevo cada semana}
\newcommand{\lblSubtituloPortadaMetodo}{Cuaderno diario para contar, trazar los números y empezar a sumar -- con la familia de \emph{Aprendo a leer}}
\newcommand{\lblPortadaNombre}{Nombre}
\newcommand{\lblPortadaCurso}{Curso}

% La medalla del final de cada trimestre (T1-T3) y el cartel de cada
% diez páginas.
\newcommand{\lblMedalla}[1]{¡Medalla de #1!}
\newcommand{\lblMedallaAnimo}{¡Sigue así, \rule{55mm}{0.4pt}!}
\newcommand{\lblPaginasHechas}[1]{¡#1 páginas hechas!}

% La clave de respuestas.
\newcommand{\lblClave}{Clave de respuestas}


% =========================================================================
% Paleta
% =========================================================================
% La de "Aprendo a leer": cada actividad tiene un color saturado (borde y
% título) y uno claro (fondo), y el color es un canal EXTRA, nunca el
% único -- el título en negrita ya dice qué actividad es. main-bw.tex fija
% \bookcolor{bw} antes de % preamble.tex
% -----------------------------------------------------------------------
% "Aprendo los números" -- cuaderno diario de números (español, A4).
% Todo el aparato tipográfico vive aquí: paquetes, colores, cajas y las
% macros que usa cada página. Es el motor de "Aprendo a leer"
% (https://github.com/konradcinkusz/learning-to-read) -- la misma letra,
% la misma paleta, la misma cabecera de cada día, las mismas cajas --,
% con lo que cambia al contar en vez de leer: la caja de arriba es "El
% número de hoy" (el número, su nombre, el marco de diez y las cosas que
% se cuentan), y las actividades son las de contar (ver notes/01-plan.md).
%
% Requiere LuaLaTeX (fontspec): la letra es Andika (SIL, OFL,
% fonts/andika/), pensada para primeros lectores -- y sus números son los
% mismos que se trazan en la actividad "Traza" (content/numeros-trazo.json
% sale del mismo glifo, ver tools/gen_numeros_puntos.py).
% -----------------------------------------------------------------------

\usepackage{fontspec}
\setmainfont{Andika}[
  Path = fonts/andika/,
  Extension = .ttf,
  UprightFont = *-Regular,
  BoldFont = *-Bold,
  ItalicFont = *-Italic,
  BoldItalicFont = *-BoldItalic,
]
\setsansfont{Andika}[
  Path = fonts/andika/,
  Extension = .ttf,
  UprightFont = *-Regular,
  BoldFont = *-Bold,
  ItalicFont = *-Italic,
  BoldItalicFont = *-BoldItalic,
]

\usepackage[spanish]{babel}
% babel-spanish hace de " un carácter activo; las comillas escritas a mano
% en el texto se estropean sin esto (ver el mismo comentario, más largo,
% en el preamble.tex de "Aprendo a leer"). Solo funciona así, al empezar
% el documento.
\AtBeginDocument{\shorthandoff{"}}

% El pie de página ("Día N / 260") vive en el pie de verdad, no en la
% página de texto: la caja de actividad llena todo lo que queda de página
% (ver `cajadia`), y un pie dentro del texto la empujaría a otra página.
\usepackage[a4paper,top=18mm,bottom=22mm,left=22mm,right=18mm,footskip=10mm]{geometry}

\usepackage{xcolor}
\usepackage[most]{tcolorbox}
\usepackage{tabularx}
\usepackage{booktabs}
\usepackage{enumitem}
\usepackage{tikz}
\usetikzlibrary{arrows.meta}
\usepackage{graphicx}
\usepackage{multicol}

\pagestyle{empty}
\setlength{\parindent}{0pt}
\setlength{\parskip}{2pt}

% --- cadenas de texto ----------------------------------------------------
% lang/es.tex
% -----------------------------------------------------------------------
% Todas las cadenas de texto del cuaderno, en un solo sitio, como en
% "Aprendo a leer": cambiar una palabra aquí la cambia en las 260 páginas
% a la vez, sin tocar el contenido generado.
% -----------------------------------------------------------------------

% Total de días del cuaderno. tools/gen_numeros.py --check comprueba que
% coincide con su propio TOTAL_DIAS.
\newcommand{\totaldias}{260}

% La cabecera de cada día (\cabeceraDia, preamble.tex).
\newcommand{\lblFecha}{Fecha}
\newcommand{\lblNotas}{Notas}
\newcommand{\lblDia}{Día}
\newcommand{\lblSemana}{Semana}
\newcommand{\lblTrimestre}{Trimestre}
\newcommand{\lblHecho}{Hecho}
\newcommand{\lblHechoSola}{sola}
\newcommand{\lblHechoConAyuda}{con ayuda}
\newcommand{\lblHechoConDificultad}{con dificultad}

% La caja de arriba, la misma todos los días.
\newcommand{\lblNumeroDeHoy}{El número de hoy}

% Los títulos de las cajas de actividad.
\newcommand{\lblTraza}{Traza}
\newcommand{\lblColorea}{Colorea}
\newcommand{\lblRodea}{Rodea}
\newcommand{\lblCuenta}{Cuenta}
\newcommand{\lblBusca}{Busca}
\newcommand{\lblUne}{Une}
\newcommand{\lblCompleta}{Completa}
\newcommand{\lblDecena}{Diez y más}
\newcommand{\lblCompara}{Compara}
\newcommand{\lblVecinos}{Antes y después}
\newcommand{\lblRecta}{La recta}
\newcommand{\lblOrdena}{Ordena}
\newcommand{\lblJunta}{Junta}
\newcommand{\lblSuma}{Suma}
\newcommand{\lblResta}{Resta}
\newcommand{\lblParte}{Parte en dos}
\newcommand{\lblDiez}{Hasta el 10}
\newcommand{\lblProblema}{Problema}
\newcommand{\lblBloques}{Decenas y unidades}
% "Number names" solo sale en "First Numbers" (lang/en.tex), pero su caja
% es la misma en los dos cuadernos (preamble.tex).
\newcommand{\lblNombres}{Los nombres de los números}
\newcommand{\lblTabla}{La tabla del 100}
\newcommand{\lblDinero}{El dinero}
\newcommand{\lblHora}{La hora}
\newcommand{\lblMide}{Mide}
\newcommand{\lblDibuja}{Dibuja}
\newcommand{\lblRepasa}{Repasa}

% Las instrucciones: las lee un adulto en voz alta (van pequeñas y en
% gris, como en "Aprendo a leer"). Las de "colorea", "rodea", "cuenta",
% "une" y "completa" las compone tools/gen_numeros.py con el número y las
% cosas de cada día ("Colorea 2 pelotas.").
\newcommand{\lblInstruccionTraza}{Repasa los puntos con el dedo, y luego con un lápiz. Después, escríbelo tú en la línea.}
\newcommand{\lblInstruccionBusca}{Rodea cada uno que encuentres. ¿Cuántos hay?}

% La tabla de las decenas (D) y las unidades (U) de la caja de arriba.
\newcommand{\lblDecenasCorto}{D}
\newcommand{\lblUnidadesCorto}{U}

% La portada.
\newcommand{\lblTituloPortada}{Aprendo los números}
\newcommand{\lblSubtituloPortada}{Un número nuevo cada semana}
\newcommand{\lblSubtituloPortadaMetodo}{Cuaderno diario para contar, trazar los números y empezar a sumar -- con la familia de \emph{Aprendo a leer}}
\newcommand{\lblPortadaNombre}{Nombre}
\newcommand{\lblPortadaCurso}{Curso}

% La medalla del final de cada trimestre (T1-T3) y el cartel de cada
% diez páginas.
\newcommand{\lblMedalla}[1]{¡Medalla de #1!}
\newcommand{\lblMedallaAnimo}{¡Sigue así, \rule{55mm}{0.4pt}!}
\newcommand{\lblPaginasHechas}[1]{¡#1 páginas hechas!}

% La clave de respuestas.
\newcommand{\lblClave}{Clave de respuestas}


% =========================================================================
% Paleta
% =========================================================================
% La de "Aprendo a leer": cada actividad tiene un color saturado (borde y
% título) y uno claro (fondo), y el color es un canal EXTRA, nunca el
% único -- el título en negrita ya dice qué actividad es. main-bw.tex fija
% \bookcolor{bw} antes de % preamble.tex
% -----------------------------------------------------------------------
% "Aprendo los números" -- cuaderno diario de números (español, A4).
% Todo el aparato tipográfico vive aquí: paquetes, colores, cajas y las
% macros que usa cada página. Es el motor de "Aprendo a leer"
% (https://github.com/konradcinkusz/learning-to-read) -- la misma letra,
% la misma paleta, la misma cabecera de cada día, las mismas cajas --,
% con lo que cambia al contar en vez de leer: la caja de arriba es "El
% número de hoy" (el número, su nombre, el marco de diez y las cosas que
% se cuentan), y las actividades son las de contar (ver notes/01-plan.md).
%
% Requiere LuaLaTeX (fontspec): la letra es Andika (SIL, OFL,
% fonts/andika/), pensada para primeros lectores -- y sus números son los
% mismos que se trazan en la actividad "Traza" (content/numeros-trazo.json
% sale del mismo glifo, ver tools/gen_numeros_puntos.py).
% -----------------------------------------------------------------------

\usepackage{fontspec}
\setmainfont{Andika}[
  Path = fonts/andika/,
  Extension = .ttf,
  UprightFont = *-Regular,
  BoldFont = *-Bold,
  ItalicFont = *-Italic,
  BoldItalicFont = *-BoldItalic,
]
\setsansfont{Andika}[
  Path = fonts/andika/,
  Extension = .ttf,
  UprightFont = *-Regular,
  BoldFont = *-Bold,
  ItalicFont = *-Italic,
  BoldItalicFont = *-BoldItalic,
]

\usepackage[spanish]{babel}
% babel-spanish hace de " un carácter activo; las comillas escritas a mano
% en el texto se estropean sin esto (ver el mismo comentario, más largo,
% en el preamble.tex de "Aprendo a leer"). Solo funciona así, al empezar
% el documento.
\AtBeginDocument{\shorthandoff{"}}

% El pie de página ("Día N / 260") vive en el pie de verdad, no en la
% página de texto: la caja de actividad llena todo lo que queda de página
% (ver `cajadia`), y un pie dentro del texto la empujaría a otra página.
\usepackage[a4paper,top=18mm,bottom=22mm,left=22mm,right=18mm,footskip=10mm]{geometry}

\usepackage{xcolor}
\usepackage[most]{tcolorbox}
\usepackage{tabularx}
\usepackage{booktabs}
\usepackage{enumitem}
\usepackage{tikz}
\usetikzlibrary{arrows.meta}
\usepackage{graphicx}
\usepackage{multicol}

\pagestyle{empty}
\setlength{\parindent}{0pt}
\setlength{\parskip}{2pt}

% --- cadenas de texto ----------------------------------------------------
\input{lang/es}

% =========================================================================
% Paleta
% =========================================================================
% La de "Aprendo a leer": cada actividad tiene un color saturado (borde y
% título) y uno claro (fondo), y el color es un canal EXTRA, nunca el
% único -- el título en negrita ya dice qué actividad es. main-bw.tex fija
% \bookcolor{bw} antes de \input{preamble}, y los mismos nombres de color
% apuntan a negro y gris: la misma página, sin tinta de color.
\makeatletter
\newif\ifbwbook
\def\aln@bookcolor@bw{bw}
\@ifundefined{bookcolor}{\bwbookfalse}{%
  \ifx\bookcolor\aln@bookcolor@bw \bwbooktrue \else \bwbookfalse \fi
}
\makeatother

\ifbwbook
  \definecolor{colorLectura}{gray}{0}
  \definecolor{colorLecturaClaro}{gray}{0.95}
  \definecolor{colorDibuja}{gray}{0}
  \definecolor{colorDibujaClaro}{gray}{0.95}
  \definecolor{colorCompleta}{gray}{0}
  \definecolor{colorCompletaClaro}{gray}{0.95}
  \definecolor{colorResponde}{gray}{0}
  \definecolor{colorRespondeClaro}{gray}{0.95}
  \definecolor{colorRelaciona}{gray}{0}
  \definecolor{colorRelacionaClaro}{gray}{0.95}
  \definecolor{colorCrea}{gray}{0}
  \definecolor{colorCreaClaro}{gray}{0.95}
  \definecolor{colorGris}{gray}{0.38}
  \definecolor{colorPunto}{gray}{0.25}
\else
  \definecolor{colorLectura}{HTML}{2E7D32}
  \definecolor{colorLecturaClaro}{HTML}{EAF6EC}
  \definecolor{colorDibuja}{HTML}{1565C0}
  \definecolor{colorDibujaClaro}{HTML}{E7F0FB}
  \definecolor{colorCompleta}{HTML}{C2185B}
  \definecolor{colorCompletaClaro}{HTML}{FBE9EF}
  \definecolor{colorResponde}{HTML}{6A1B9A}
  \definecolor{colorRespondeClaro}{HTML}{F2E9F7}
  \definecolor{colorRelaciona}{HTML}{E65100}
  \definecolor{colorRelacionaClaro}{HTML}{FDEEE2}
  \definecolor{colorCrea}{HTML}{AD1457}
  \definecolor{colorCreaClaro}{HTML}{FBE7EE}
  \definecolor{colorGris}{HTML}{616161}
  \definecolor{colorPunto}{HTML}{6A1B9A}
\fi

% =========================================================================
% Los dibujos
% =========================================================================
% diagrams/kit.tex: las piezas de los dibujos de «First Words» (de
% "Aprendo a leer"), con las que están hechas muchas de las cosas que se
% cuentan; diagrams/objetos.tex: esas cosas, cada una en la misma caja.
\input{diagrams/kit}
\input{diagrams/objetos}

% =========================================================================
% Cajas
% =========================================================================
% La caja de actividad llena todo lo que queda de página (`height fill`),
% como en el cuaderno de primeras palabras: casi todas las actividades
% terminan coloreando o dibujando, y el sitio es todo el que haya. Lo que
% tiene que ir abajo del todo (la línea de escribir, el cartel de cada
% diez páginas) va en la parte de abajo de la caja (\tcblower, estilo
% `abajo`); lo que tiene que ir centrado en el hueco, con `centrado abajo`.
\tcbset{
  cajabase/.style={
    enhanced,
    boxrule=1.1pt,
    arc=3mm,
    left=6mm, right=6mm, top=5mm, bottom=5mm,
    fonttitle=\bfseries\sffamily\large,
    coltitle=white,
  },
  cajadia/.style={cajabase, height fill},
  abajo/.style={space to upper, segmentation hidden, middle=2mm},
  centrado abajo/.style={space to lower, valign lower=center, segmentation hidden, middle=2mm},
}

\newtcolorbox{cajaNumero}{cajabase, top=3mm, bottom=3mm,
  colback=colorLecturaClaro, colframe=colorLectura, colbacktitle=colorLectura,
  title=\lblNumeroDeHoy}

\newtcolorbox{cajaTraza}[1][]{cajadia, #1,
  colback=colorRespondeClaro, colframe=colorResponde, colbacktitle=colorResponde,
  title=\lblTraza}
\newtcolorbox{cajaColorea}[1][]{cajadia, #1,
  colback=colorCompletaClaro, colframe=colorCompleta, colbacktitle=colorCompleta,
  title=\lblColorea}
\newtcolorbox{cajaRodea}[1][]{cajadia, #1,
  colback=colorRelacionaClaro, colframe=colorRelaciona, colbacktitle=colorRelaciona,
  title=\lblRodea}
\newtcolorbox{cajaCuenta}[1][]{cajadia, #1,
  colback=colorRelacionaClaro, colframe=colorRelaciona, colbacktitle=colorRelaciona,
  title=\lblCuenta}
\newtcolorbox{cajaBusca}[1][]{cajadia, #1,
  colback=colorCompletaClaro, colframe=colorCompleta, colbacktitle=colorCompleta,
  title=\lblBusca}
\newtcolorbox{cajaUne}[1][]{cajadia, #1,
  colback=colorRelacionaClaro, colframe=colorRelaciona, colbacktitle=colorRelaciona,
  title=\lblUne}
\newtcolorbox{cajaCompleta}[1][]{cajadia, #1,
  colback=colorRespondeClaro, colframe=colorResponde, colbacktitle=colorResponde,
  title=\lblCompleta}
% Las del invierno.
\newtcolorbox{cajaDecena}[1][]{cajadia, #1,
  colback=colorRelacionaClaro, colframe=colorRelaciona, colbacktitle=colorRelaciona,
  title=\lblDecena}
\newtcolorbox{cajaCompara}[1][]{cajadia, #1,
  colback=colorRelacionaClaro, colframe=colorRelaciona, colbacktitle=colorRelaciona,
  title=\lblCompara}
\newtcolorbox{cajaVecinos}[1][]{cajadia, #1,
  colback=colorRespondeClaro, colframe=colorResponde, colbacktitle=colorResponde,
  title=\lblVecinos}
\newtcolorbox{cajaRecta}[1][]{cajadia, #1,
  colback=colorRespondeClaro, colframe=colorResponde, colbacktitle=colorResponde,
  title=\lblRecta}
\newtcolorbox{cajaOrdena}[1][]{cajadia, #1,
  colback=colorCompletaClaro, colframe=colorCompleta, colbacktitle=colorCompleta,
  title=\lblOrdena}
\newtcolorbox{cajaJunta}[1][]{cajadia, #1,
  colback=colorDibujaClaro, colframe=colorDibuja, colbacktitle=colorDibuja,
  title=\lblJunta}
\newtcolorbox{cajaDibuja}[1][]{cajadia, #1,
  colback=colorDibujaClaro, colframe=colorDibuja, colbacktitle=colorDibuja,
  title=\lblDibuja}
\newtcolorbox{cajaRepasa}[1][]{cajadia, #1,
  colback=colorCreaClaro, colframe=colorCrea, colbacktitle=colorCrea,
  title=\lblRepasa}

% La instrucción de cada actividad: la lee el adulto.
\newcommand{\instruccion}[1]{{\footnotesize\color{colorGris}#1\par}}
% El enunciado grande, lo que se hace (lo que el niño o la niña puede
% seguir con el dedo): "Colorea 2 pelotas."
\newcommand{\enunciado}[1]{{\LARGE #1\par}}
% Un renglón para escribir el número: el número de muestra, en gris, y
% sitio en la misma línea para escribirlo tres o cuatro veces más.
\newcommand{\renglonTraza}[1]{%
  \par\noindent\begin{tikzpicture}
    \draw[line width=0.6pt] (0,0) -- (\linewidth-1pt,0);
    \node[anchor=base west, inner sep=0pt, text=colorGris!55,
      font=\fontsize{44}{44}\selectfont\bfseries] at (2mm,0) {#1};
  \end{tikzpicture}\par\vspace{12mm}}
% El hueco en el que se escribe la respuesta: "10 y 3 son [ ]".
\newcommand{\huecoRespuesta}{\raisebox{-5mm}{\tikz\draw[line width=1.1pt,
  rounded corners=2mm, fill=white] (0,0) rectangle (26mm,18mm);}}
% Lo que no cabe en su sitio, encogido hasta que quepa (#1, el ancho):
% los nombres largos, como "diecisiete", en la caja de arriba.
\newsavebox{\cajaAjustada}
\newcommand{\ajustado}[2]{\sbox{\cajaAjustada}{#2}%
  \ifdim\wd\cajaAjustada>#1\relax\resizebox{#1}{!}{\usebox{\cajaAjustada}}%
  \else\usebox{\cajaAjustada}\fi}
% El cartel de cada diez páginas, al pie de "Repasa".
\newcommand{\cartel}[1]{\begin{center}\bfseries\Large\color{colorCrea}#1\end{center}}

\newcommand{\casilla}{\raisebox{-0.15ex}{\framebox[3.4mm][c]{\rule{0pt}{2.6mm}}}}
\newenvironment{listaRepaso}{%
  \begin{itemize}[label=\casilla, leftmargin=8mm, itemsep=2mm, topsep=1mm]\large
}{%
  \end{itemize}
}

% =========================================================================
% La página de un día
% =========================================================================
% La cabecera es la de "Aprendo a leer": el día, la semana y el
% trimestre; el tema de la semana (el mismo que esa semana en "Aprendo a
% leer", para que quien lleve los dos cuadernos cuente lo que lee), y los
% campos para el adulto -- aquí "Hecho" en vez de "Leído". Deja
% \label{dia:N}: tools/check_pages.py comprueba con esas etiquetas que
% cada día ocupa exactamente una página.
\newcommand{\cabeceraDia}[4]{%
  \label{dia:#1}%
  \noindent{\LARGE\bfseries\lblDia~#1}\quad
  {\large(\lblSemana~#2 \textperiodcentered\ \lblTrimestre~#3)}\par
  {\normalsize\bfseries\color{colorGris}#4}\par
  \vspace{1mm}
  \textbf{\lblFecha:} \rule{22mm}{0.4pt}\hspace{4mm}%
  \textbf{\lblHecho:} \casilla\ \lblHechoSola\quad
  \casilla\ \lblHechoConAyuda\quad
  \casilla\ \lblHechoConDificultad\hspace{4mm}%
  \textbf{\lblNotas:} \rule{20mm}{0.4pt}\par
  \vspace{3mm}
}

\makeatletter
\def\ps@pienumeros{%
  \let\@oddhead\@empty\let\@evenhead\@empty
  \def\@oddfoot{\hfil\footnotesize\color{colorGris}\lblDia~\diapagina@actual~/ \totaldias\hfil}%
  \let\@evenfoot\@oddfoot
}
\newenvironment{diapagina}[4]{%
  \gdef\diapagina@actual{#1}%
  \thispagestyle{pienumeros}%
  \cabeceraDia{#1}{#2}{#3}{#4}%
}{%
  \par\newpage
}

% Una etiqueta que solo guarda la página en la que cae, para
% tools/check_pages.py: la clave de respuestas empieza con
% \etiquetaPagina{clave}, así se comprueba también que el último día no
% se ha desbordado a una segunda página.
\newcommand{\etiquetaPagina}[1]{\begingroup\def\@currentlabel{}\label{#1}\endgroup}
\makeatother

% =========================================================================
% El número de hoy
% =========================================================================
% El marco de diez: dos filas de cinco casillas, con un punto en las N
% primeras -- la manera de ver de un vistazo cuánto es un número hasta
% el 10 (5 es una fila entera; 7, una fila y dos más).
\newcommand{\marcoDiez}[1]{%
  \begin{tikzpicture}[x=7.5mm, y=7.5mm, line width=0.9pt]
    \foreach \i in {0,...,9} {
      \pgfmathtruncatemacro\col{mod(\i,5)}
      \pgfmathtruncatemacro\fila{\i<5 ? 1 : 0}
      \draw (\col,\fila) rectangle ++(1,1);
      \ifnum\i<#1\relax
        \fill[colorLectura] (\col+0.5,\fila+0.5) circle (0.32);
      \fi
    }
  \end{tikzpicture}}
% Del 11 al 20, dos marcos, uno encima de otro: el primero, lleno.
\newcommand{\marcosDiez}[1]{%
  \ifnum#1>10\relax
    \marcoDiez{10}\par\vspace{1.5mm}\marcoDiez{\numexpr#1-10\relax}%
  \else
    \marcoDiez{#1}%
  \fi}

% La caja de arriba: el número grande, su nombre, el marco de diez y las
% cosas que se cuentan ese día (#3: un tikzpicture que compone
% tools/gen_numeros.py), y debajo la frase de la historia, que lee el
% adulto. La columna del número mide 28 mm, o lo que mida el número si
% tiene dos cifras (el 10): no se encoge, y lo que se estrecha es la de
% las cosas (tools/gen_numeros.py las dibuja más pequeñas).
\newlength{\anchoNumeroDeHoy}
\newcommand{\fuenteNumeroDeHoy}{\fontsize{92}{92}\selectfont\bfseries}
\newcommand{\numeroDeHoy}[4]{%
  \settowidth{\anchoNumeroDeHoy}{\fuenteNumeroDeHoy #1}%
  \ifdim\anchoNumeroDeHoy<28mm \setlength{\anchoNumeroDeHoy}{28mm}\fi
  \begin{cajaNumero}
  \begin{tabularx}{\linewidth}{@{}>{\centering\arraybackslash}m{\anchoNumeroDeHoy}@{\hspace{3mm}}>{\centering\arraybackslash}m{42mm}@{\hspace{3mm}}>{\centering\arraybackslash}X@{}}
  {\fuenteNumeroDeHoy\color{colorLectura}#1} &
  \ajustado{42mm}{\fontsize{30}{34}\selectfont\bfseries #2}\par\vspace{3mm}\marcosDiez{#1} &
  #3 \\
  \end{tabularx}
  \par\vspace{2mm}
  {\large #4\par}
  \end{cajaNumero}\par\vspace{4mm}}

% =========================================================================
% Medalla de fin de trimestre (tools/gen_numeros.py, PLANTILLA_MEDALLA)
% =========================================================================
\newcommand{\medalla}[1]{%
  \begin{tikzpicture}[line width=1.6pt, color=colorCrea]
    \draw (0,0) circle (1.1);
    \draw (0,0) circle (0.85);
    \node at (0,0) {\Large\bfseries #1};
    \draw (-0.35,-1.05) -- (-0.7,-2.2) -- (-0.15,-1.75) -- cycle;
    \draw (0.35,-1.05)  -- (0.7,-2.2)  -- (0.15,-1.75)  -- cycle;
  \end{tikzpicture}}

% =========================================================================
% La clave de respuestas
% =========================================================================
\newcommand{\claveEntrada}[3]{%
  \item \textbf{\lblDia\ #1} \textcolor{colorGris}{(#2)}: #3%
}
, y los mismos nombres de color
% apuntan a negro y gris: la misma página, sin tinta de color.
\makeatletter
\newif\ifbwbook
\def\aln@bookcolor@bw{bw}
\@ifundefined{bookcolor}{\bwbookfalse}{%
  \ifx\bookcolor\aln@bookcolor@bw \bwbooktrue \else \bwbookfalse \fi
}
\makeatother

\ifbwbook
  \definecolor{colorLectura}{gray}{0}
  \definecolor{colorLecturaClaro}{gray}{0.95}
  \definecolor{colorDibuja}{gray}{0}
  \definecolor{colorDibujaClaro}{gray}{0.95}
  \definecolor{colorCompleta}{gray}{0}
  \definecolor{colorCompletaClaro}{gray}{0.95}
  \definecolor{colorResponde}{gray}{0}
  \definecolor{colorRespondeClaro}{gray}{0.95}
  \definecolor{colorRelaciona}{gray}{0}
  \definecolor{colorRelacionaClaro}{gray}{0.95}
  \definecolor{colorCrea}{gray}{0}
  \definecolor{colorCreaClaro}{gray}{0.95}
  \definecolor{colorGris}{gray}{0.38}
  \definecolor{colorPunto}{gray}{0.25}
\else
  \definecolor{colorLectura}{HTML}{2E7D32}
  \definecolor{colorLecturaClaro}{HTML}{EAF6EC}
  \definecolor{colorDibuja}{HTML}{1565C0}
  \definecolor{colorDibujaClaro}{HTML}{E7F0FB}
  \definecolor{colorCompleta}{HTML}{C2185B}
  \definecolor{colorCompletaClaro}{HTML}{FBE9EF}
  \definecolor{colorResponde}{HTML}{6A1B9A}
  \definecolor{colorRespondeClaro}{HTML}{F2E9F7}
  \definecolor{colorRelaciona}{HTML}{E65100}
  \definecolor{colorRelacionaClaro}{HTML}{FDEEE2}
  \definecolor{colorCrea}{HTML}{AD1457}
  \definecolor{colorCreaClaro}{HTML}{FBE7EE}
  \definecolor{colorGris}{HTML}{616161}
  \definecolor{colorPunto}{HTML}{6A1B9A}
\fi

% =========================================================================
% Los dibujos
% =========================================================================
% diagrams/kit.tex: las piezas de los dibujos de «First Words» (de
% "Aprendo a leer"), con las que están hechas muchas de las cosas que se
% cuentan; diagrams/objetos.tex: esas cosas, cada una en la misma caja.
\input{diagrams/kit}
% diagrams/objetos.tex -- las cosas que se cuentan: cada una dibujada
% dentro de la misma caja, de -1 a 1 en x y en y, para que
% tools/gen_numeros.py las pueda poner en fila, en grupos o en el marco
% de diez sin saber qué es cada una (ver OBJETOS en tools/gen_numeros.py,
% que dice qué macro dibuja cada cosa y cómo se llama en singular y en
% plural). Todas las zonas quedan cerradas y rellenas de blanco, para
% que cada una se pueda colorear.
%
% Muchas son los dibujos de «First Words» (diagrams/kit.tex, del
% repositorio de "Aprendo a leer"), encogidos o movidos hasta caber en la
% caja; las demás son nuevas y siguen el mismo estilo de línea.
\tikzset{objeto/.style={line width=1.1pt, line cap=round, line join=round}}

% Una manzana, con el rabito y una hoja.
\newcommand{\objManzana}{%
  \filldraw[fill=white] (0,0.55) .. controls (-0.35,0.9) and (-1.0,0.78) .. (-0.95,0.02)
    .. controls (-0.9,-0.72) and (-0.35,-1.0) .. (0,-0.8)
    .. controls (0.35,-1.0) and (0.9,-0.72) .. (0.95,0.02)
    .. controls (1.0,0.78) and (0.35,0.9) .. (0,0.55) -- cycle;
  \draw (0,0.55) .. controls (0.02,0.8) .. (0.1,0.98);
  \filldraw[fill=white] (0.08,0.85) .. controls (0.32,1.05) and (0.62,0.98) .. (0.7,0.84)
    .. controls (0.5,0.68) and (0.25,0.7) .. (0.08,0.85) -- cycle;}

% Una pelota de playa, con cuatro gajos.
\newcommand{\objPelota}{%
  \filldraw[fill=white] (0,0) circle (0.9);
  \draw (-0.9,0) .. controls (-0.3,0.38) and (0.3,0.38) .. (0.9,0);
  \draw (0,0.9) .. controls (-0.38,0.3) and (-0.38,-0.3) .. (0,-0.9);}

% Un hueso, tumbado: la barra y los cuatro nudos, un solo contorno.
\newcommand{\objHueso}{%
  \filldraw[fill=white] (-0.492,0.16) -- (0.492,0.16)
    arc[start angle=197.2, end angle=-62.7, radius=0.27]
    arc[start angle=62.7, end angle=-197.2, radius=0.27]
    -- (-0.492,-0.16)
    arc[start angle=17.2, end angle=-242.7, radius=0.27]
    arc[start angle=242.7, end angle=-17.2, radius=0.27] -- cycle;}

% Un sol: el círculo y ocho rayos.
\newcommand{\objSol}{%
  \filldraw[fill=white] (0,0) circle (0.5);
  \foreach \a in {0,45,...,315} {\draw (\a:0.66) -- (\a:0.96);}}

% Una huella de Toby.
\newcommand{\objHuella}{%
  \begin{scope}[shift={(0,-0.28)}, scale=1.65] \huellaFW \end{scope}}

% Toby, sentado.
\newcommand{\objToby}{%
  \begin{scope}[shift={(0.38,-0.88)}, scale=0.4] \tobySentadoFW \end{scope}}

% Un globo, con su hilo.
\newcommand{\objGlobo}{%
  \begin{scope}[shift={(0,-0.85)}, scale=0.72] \globoFW{0.2}{-0.2} \end{scope}}

% Una estrella.
\newcommand{\objEstrella}{%
  \begin{scope}[scale=0.95] \estrellaFW \end{scope}}

% Una hoja de otoño.
\newcommand{\objHoja}{%
  \begin{scope}[shift={(0,-0.8)}, scale=1.2] \hojaFW \end{scope}}

% Un corazón.
\newcommand{\objCorazon}{%
  \begin{scope}[shift={(0,-0.82)}, scale=1.35] \corazonFW \end{scope}}

% Un caramelo.
\newcommand{\objCaramelo}{%
  \begin{scope}[scale=1.25] \carameloFW \end{scope}}

% Un pez.
\newcommand{\objPez}{%
  \begin{scope}[shift={(0.2,0)}, scale=0.72] \pezFW \end{scope}}


% =========================================================================
% Cajas
% =========================================================================
% La caja de actividad llena todo lo que queda de página (`height fill`),
% como en el cuaderno de primeras palabras: casi todas las actividades
% terminan coloreando o dibujando, y el sitio es todo el que haya. Lo que
% tiene que ir abajo del todo (la línea de escribir, el cartel de cada
% diez páginas) va en la parte de abajo de la caja (\tcblower, estilo
% `abajo`); lo que tiene que ir centrado en el hueco, con `centrado abajo`.
\tcbset{
  cajabase/.style={
    enhanced,
    boxrule=1.1pt,
    arc=3mm,
    left=6mm, right=6mm, top=5mm, bottom=5mm,
    fonttitle=\bfseries\sffamily\large,
    coltitle=white,
  },
  cajadia/.style={cajabase, height fill},
  abajo/.style={space to upper, segmentation hidden, middle=2mm},
  centrado abajo/.style={space to lower, valign lower=center, segmentation hidden, middle=2mm},
}

\newtcolorbox{cajaNumero}{cajabase, top=3mm, bottom=3mm,
  colback=colorLecturaClaro, colframe=colorLectura, colbacktitle=colorLectura,
  title=\lblNumeroDeHoy}

\newtcolorbox{cajaTraza}[1][]{cajadia, #1,
  colback=colorRespondeClaro, colframe=colorResponde, colbacktitle=colorResponde,
  title=\lblTraza}
\newtcolorbox{cajaColorea}[1][]{cajadia, #1,
  colback=colorCompletaClaro, colframe=colorCompleta, colbacktitle=colorCompleta,
  title=\lblColorea}
\newtcolorbox{cajaRodea}[1][]{cajadia, #1,
  colback=colorRelacionaClaro, colframe=colorRelaciona, colbacktitle=colorRelaciona,
  title=\lblRodea}
\newtcolorbox{cajaCuenta}[1][]{cajadia, #1,
  colback=colorRelacionaClaro, colframe=colorRelaciona, colbacktitle=colorRelaciona,
  title=\lblCuenta}
\newtcolorbox{cajaBusca}[1][]{cajadia, #1,
  colback=colorCompletaClaro, colframe=colorCompleta, colbacktitle=colorCompleta,
  title=\lblBusca}
\newtcolorbox{cajaUne}[1][]{cajadia, #1,
  colback=colorRelacionaClaro, colframe=colorRelaciona, colbacktitle=colorRelaciona,
  title=\lblUne}
\newtcolorbox{cajaCompleta}[1][]{cajadia, #1,
  colback=colorRespondeClaro, colframe=colorResponde, colbacktitle=colorResponde,
  title=\lblCompleta}
% Las del invierno.
\newtcolorbox{cajaDecena}[1][]{cajadia, #1,
  colback=colorRelacionaClaro, colframe=colorRelaciona, colbacktitle=colorRelaciona,
  title=\lblDecena}
\newtcolorbox{cajaCompara}[1][]{cajadia, #1,
  colback=colorRelacionaClaro, colframe=colorRelaciona, colbacktitle=colorRelaciona,
  title=\lblCompara}
\newtcolorbox{cajaVecinos}[1][]{cajadia, #1,
  colback=colorRespondeClaro, colframe=colorResponde, colbacktitle=colorResponde,
  title=\lblVecinos}
\newtcolorbox{cajaRecta}[1][]{cajadia, #1,
  colback=colorRespondeClaro, colframe=colorResponde, colbacktitle=colorResponde,
  title=\lblRecta}
\newtcolorbox{cajaOrdena}[1][]{cajadia, #1,
  colback=colorCompletaClaro, colframe=colorCompleta, colbacktitle=colorCompleta,
  title=\lblOrdena}
\newtcolorbox{cajaJunta}[1][]{cajadia, #1,
  colback=colorDibujaClaro, colframe=colorDibuja, colbacktitle=colorDibuja,
  title=\lblJunta}
\newtcolorbox{cajaDibuja}[1][]{cajadia, #1,
  colback=colorDibujaClaro, colframe=colorDibuja, colbacktitle=colorDibuja,
  title=\lblDibuja}
\newtcolorbox{cajaRepasa}[1][]{cajadia, #1,
  colback=colorCreaClaro, colframe=colorCrea, colbacktitle=colorCrea,
  title=\lblRepasa}

% La instrucción de cada actividad: la lee el adulto.
\newcommand{\instruccion}[1]{{\footnotesize\color{colorGris}#1\par}}
% El enunciado grande, lo que se hace (lo que el niño o la niña puede
% seguir con el dedo): "Colorea 2 pelotas."
\newcommand{\enunciado}[1]{{\LARGE #1\par}}
% Un renglón para escribir el número: el número de muestra, en gris, y
% sitio en la misma línea para escribirlo tres o cuatro veces más.
\newcommand{\renglonTraza}[1]{%
  \par\noindent\begin{tikzpicture}
    \draw[line width=0.6pt] (0,0) -- (\linewidth-1pt,0);
    \node[anchor=base west, inner sep=0pt, text=colorGris!55,
      font=\fontsize{44}{44}\selectfont\bfseries] at (2mm,0) {#1};
  \end{tikzpicture}\par\vspace{12mm}}
% El hueco en el que se escribe la respuesta: "10 y 3 son [ ]".
\newcommand{\huecoRespuesta}{\raisebox{-5mm}{\tikz\draw[line width=1.1pt,
  rounded corners=2mm, fill=white] (0,0) rectangle (26mm,18mm);}}
% Lo que no cabe en su sitio, encogido hasta que quepa (#1, el ancho):
% los nombres largos, como "diecisiete", en la caja de arriba.
\newsavebox{\cajaAjustada}
\newcommand{\ajustado}[2]{\sbox{\cajaAjustada}{#2}%
  \ifdim\wd\cajaAjustada>#1\relax\resizebox{#1}{!}{\usebox{\cajaAjustada}}%
  \else\usebox{\cajaAjustada}\fi}
% El cartel de cada diez páginas, al pie de "Repasa".
\newcommand{\cartel}[1]{\begin{center}\bfseries\Large\color{colorCrea}#1\end{center}}

\newcommand{\casilla}{\raisebox{-0.15ex}{\framebox[3.4mm][c]{\rule{0pt}{2.6mm}}}}
\newenvironment{listaRepaso}{%
  \begin{itemize}[label=\casilla, leftmargin=8mm, itemsep=2mm, topsep=1mm]\large
}{%
  \end{itemize}
}

% =========================================================================
% La página de un día
% =========================================================================
% La cabecera es la de "Aprendo a leer": el día, la semana y el
% trimestre; el tema de la semana (el mismo que esa semana en "Aprendo a
% leer", para que quien lleve los dos cuadernos cuente lo que lee), y los
% campos para el adulto -- aquí "Hecho" en vez de "Leído". Deja
% \label{dia:N}: tools/check_pages.py comprueba con esas etiquetas que
% cada día ocupa exactamente una página.
\newcommand{\cabeceraDia}[4]{%
  \label{dia:#1}%
  \noindent{\LARGE\bfseries\lblDia~#1}\quad
  {\large(\lblSemana~#2 \textperiodcentered\ \lblTrimestre~#3)}\par
  {\normalsize\bfseries\color{colorGris}#4}\par
  \vspace{1mm}
  \textbf{\lblFecha:} \rule{22mm}{0.4pt}\hspace{4mm}%
  \textbf{\lblHecho:} \casilla\ \lblHechoSola\quad
  \casilla\ \lblHechoConAyuda\quad
  \casilla\ \lblHechoConDificultad\hspace{4mm}%
  \textbf{\lblNotas:} \rule{20mm}{0.4pt}\par
  \vspace{3mm}
}

\makeatletter
\def\ps@pienumeros{%
  \let\@oddhead\@empty\let\@evenhead\@empty
  \def\@oddfoot{\hfil\footnotesize\color{colorGris}\lblDia~\diapagina@actual~/ \totaldias\hfil}%
  \let\@evenfoot\@oddfoot
}
\newenvironment{diapagina}[4]{%
  \gdef\diapagina@actual{#1}%
  \thispagestyle{pienumeros}%
  \cabeceraDia{#1}{#2}{#3}{#4}%
}{%
  \par\newpage
}

% Una etiqueta que solo guarda la página en la que cae, para
% tools/check_pages.py: la clave de respuestas empieza con
% \etiquetaPagina{clave}, así se comprueba también que el último día no
% se ha desbordado a una segunda página.
\newcommand{\etiquetaPagina}[1]{\begingroup\def\@currentlabel{}\label{#1}\endgroup}
\makeatother

% =========================================================================
% El número de hoy
% =========================================================================
% El marco de diez: dos filas de cinco casillas, con un punto en las N
% primeras -- la manera de ver de un vistazo cuánto es un número hasta
% el 10 (5 es una fila entera; 7, una fila y dos más).
\newcommand{\marcoDiez}[1]{%
  \begin{tikzpicture}[x=7.5mm, y=7.5mm, line width=0.9pt]
    \foreach \i in {0,...,9} {
      \pgfmathtruncatemacro\col{mod(\i,5)}
      \pgfmathtruncatemacro\fila{\i<5 ? 1 : 0}
      \draw (\col,\fila) rectangle ++(1,1);
      \ifnum\i<#1\relax
        \fill[colorLectura] (\col+0.5,\fila+0.5) circle (0.32);
      \fi
    }
  \end{tikzpicture}}
% Del 11 al 20, dos marcos, uno encima de otro: el primero, lleno.
\newcommand{\marcosDiez}[1]{%
  \ifnum#1>10\relax
    \marcoDiez{10}\par\vspace{1.5mm}\marcoDiez{\numexpr#1-10\relax}%
  \else
    \marcoDiez{#1}%
  \fi}

% La caja de arriba: el número grande, su nombre, el marco de diez y las
% cosas que se cuentan ese día (#3: un tikzpicture que compone
% tools/gen_numeros.py), y debajo la frase de la historia, que lee el
% adulto. La columna del número mide 28 mm, o lo que mida el número si
% tiene dos cifras (el 10): no se encoge, y lo que se estrecha es la de
% las cosas (tools/gen_numeros.py las dibuja más pequeñas).
\newlength{\anchoNumeroDeHoy}
\newcommand{\fuenteNumeroDeHoy}{\fontsize{92}{92}\selectfont\bfseries}
\newcommand{\numeroDeHoy}[4]{%
  \settowidth{\anchoNumeroDeHoy}{\fuenteNumeroDeHoy #1}%
  \ifdim\anchoNumeroDeHoy<28mm \setlength{\anchoNumeroDeHoy}{28mm}\fi
  \begin{cajaNumero}
  \begin{tabularx}{\linewidth}{@{}>{\centering\arraybackslash}m{\anchoNumeroDeHoy}@{\hspace{3mm}}>{\centering\arraybackslash}m{42mm}@{\hspace{3mm}}>{\centering\arraybackslash}X@{}}
  {\fuenteNumeroDeHoy\color{colorLectura}#1} &
  \ajustado{42mm}{\fontsize{30}{34}\selectfont\bfseries #2}\par\vspace{3mm}\marcosDiez{#1} &
  #3 \\
  \end{tabularx}
  \par\vspace{2mm}
  {\large #4\par}
  \end{cajaNumero}\par\vspace{4mm}}

% =========================================================================
% Medalla de fin de trimestre (tools/gen_numeros.py, PLANTILLA_MEDALLA)
% =========================================================================
\newcommand{\medalla}[1]{%
  \begin{tikzpicture}[line width=1.6pt, color=colorCrea]
    \draw (0,0) circle (1.1);
    \draw (0,0) circle (0.85);
    \node at (0,0) {\Large\bfseries #1};
    \draw (-0.35,-1.05) -- (-0.7,-2.2) -- (-0.15,-1.75) -- cycle;
    \draw (0.35,-1.05)  -- (0.7,-2.2)  -- (0.15,-1.75)  -- cycle;
  \end{tikzpicture}}

% =========================================================================
% La clave de respuestas
% =========================================================================
\newcommand{\claveEntrada}[3]{%
  \item \textbf{\lblDia\ #1} \textcolor{colorGris}{(#2)}: #3%
}
, y los mismos nombres de color
% apuntan a negro y gris: la misma página, sin tinta de color.
\makeatletter
\newif\ifbwbook
\def\aln@bookcolor@bw{bw}
\@ifundefined{bookcolor}{\bwbookfalse}{%
  \ifx\bookcolor\aln@bookcolor@bw \bwbooktrue \else \bwbookfalse \fi
}
\makeatother

\ifbwbook
  \definecolor{colorLectura}{gray}{0}
  \definecolor{colorLecturaClaro}{gray}{0.95}
  \definecolor{colorDibuja}{gray}{0}
  \definecolor{colorDibujaClaro}{gray}{0.95}
  \definecolor{colorCompleta}{gray}{0}
  \definecolor{colorCompletaClaro}{gray}{0.95}
  \definecolor{colorResponde}{gray}{0}
  \definecolor{colorRespondeClaro}{gray}{0.95}
  \definecolor{colorRelaciona}{gray}{0}
  \definecolor{colorRelacionaClaro}{gray}{0.95}
  \definecolor{colorCrea}{gray}{0}
  \definecolor{colorCreaClaro}{gray}{0.95}
  \definecolor{colorGris}{gray}{0.38}
  \definecolor{colorPunto}{gray}{0.25}
\else
  \definecolor{colorLectura}{HTML}{2E7D32}
  \definecolor{colorLecturaClaro}{HTML}{EAF6EC}
  \definecolor{colorDibuja}{HTML}{1565C0}
  \definecolor{colorDibujaClaro}{HTML}{E7F0FB}
  \definecolor{colorCompleta}{HTML}{C2185B}
  \definecolor{colorCompletaClaro}{HTML}{FBE9EF}
  \definecolor{colorResponde}{HTML}{6A1B9A}
  \definecolor{colorRespondeClaro}{HTML}{F2E9F7}
  \definecolor{colorRelaciona}{HTML}{E65100}
  \definecolor{colorRelacionaClaro}{HTML}{FDEEE2}
  \definecolor{colorCrea}{HTML}{AD1457}
  \definecolor{colorCreaClaro}{HTML}{FBE7EE}
  \definecolor{colorGris}{HTML}{616161}
  \definecolor{colorPunto}{HTML}{6A1B9A}
\fi

% =========================================================================
% Los dibujos
% =========================================================================
% diagrams/kit.tex: las piezas de los dibujos de «First Words» (de
% "Aprendo a leer"), con las que están hechas muchas de las cosas que se
% cuentan; diagrams/objetos.tex: esas cosas, cada una en la misma caja.
\input{diagrams/kit}
% diagrams/objetos.tex -- las cosas que se cuentan: cada una dibujada
% dentro de la misma caja, de -1 a 1 en x y en y, para que
% tools/gen_numeros.py las pueda poner en fila, en grupos o en el marco
% de diez sin saber qué es cada una (ver OBJETOS en tools/gen_numeros.py,
% que dice qué macro dibuja cada cosa y cómo se llama en singular y en
% plural). Todas las zonas quedan cerradas y rellenas de blanco, para
% que cada una se pueda colorear.
%
% Muchas son los dibujos de «First Words» (diagrams/kit.tex, del
% repositorio de "Aprendo a leer"), encogidos o movidos hasta caber en la
% caja; las demás son nuevas y siguen el mismo estilo de línea.
\tikzset{objeto/.style={line width=1.1pt, line cap=round, line join=round}}

% Una manzana, con el rabito y una hoja.
\newcommand{\objManzana}{%
  \filldraw[fill=white] (0,0.55) .. controls (-0.35,0.9) and (-1.0,0.78) .. (-0.95,0.02)
    .. controls (-0.9,-0.72) and (-0.35,-1.0) .. (0,-0.8)
    .. controls (0.35,-1.0) and (0.9,-0.72) .. (0.95,0.02)
    .. controls (1.0,0.78) and (0.35,0.9) .. (0,0.55) -- cycle;
  \draw (0,0.55) .. controls (0.02,0.8) .. (0.1,0.98);
  \filldraw[fill=white] (0.08,0.85) .. controls (0.32,1.05) and (0.62,0.98) .. (0.7,0.84)
    .. controls (0.5,0.68) and (0.25,0.7) .. (0.08,0.85) -- cycle;}

% Una pelota de playa, con cuatro gajos.
\newcommand{\objPelota}{%
  \filldraw[fill=white] (0,0) circle (0.9);
  \draw (-0.9,0) .. controls (-0.3,0.38) and (0.3,0.38) .. (0.9,0);
  \draw (0,0.9) .. controls (-0.38,0.3) and (-0.38,-0.3) .. (0,-0.9);}

% Un hueso, tumbado: la barra y los cuatro nudos, un solo contorno.
\newcommand{\objHueso}{%
  \filldraw[fill=white] (-0.492,0.16) -- (0.492,0.16)
    arc[start angle=197.2, end angle=-62.7, radius=0.27]
    arc[start angle=62.7, end angle=-197.2, radius=0.27]
    -- (-0.492,-0.16)
    arc[start angle=17.2, end angle=-242.7, radius=0.27]
    arc[start angle=242.7, end angle=-17.2, radius=0.27] -- cycle;}

% Un sol: el círculo y ocho rayos.
\newcommand{\objSol}{%
  \filldraw[fill=white] (0,0) circle (0.5);
  \foreach \a in {0,45,...,315} {\draw (\a:0.66) -- (\a:0.96);}}

% Una huella de Toby.
\newcommand{\objHuella}{%
  \begin{scope}[shift={(0,-0.28)}, scale=1.65] \huellaFW \end{scope}}

% Toby, sentado.
\newcommand{\objToby}{%
  \begin{scope}[shift={(0.38,-0.88)}, scale=0.4] \tobySentadoFW \end{scope}}

% Un globo, con su hilo.
\newcommand{\objGlobo}{%
  \begin{scope}[shift={(0,-0.85)}, scale=0.72] \globoFW{0.2}{-0.2} \end{scope}}

% Una estrella.
\newcommand{\objEstrella}{%
  \begin{scope}[scale=0.95] \estrellaFW \end{scope}}

% Una hoja de otoño.
\newcommand{\objHoja}{%
  \begin{scope}[shift={(0,-0.8)}, scale=1.2] \hojaFW \end{scope}}

% Un corazón.
\newcommand{\objCorazon}{%
  \begin{scope}[shift={(0,-0.82)}, scale=1.35] \corazonFW \end{scope}}

% Un caramelo.
\newcommand{\objCaramelo}{%
  \begin{scope}[scale=1.25] \carameloFW \end{scope}}

% Un pez.
\newcommand{\objPez}{%
  \begin{scope}[shift={(0.2,0)}, scale=0.72] \pezFW \end{scope}}


% =========================================================================
% Cajas
% =========================================================================
% La caja de actividad llena todo lo que queda de página (`height fill`),
% como en el cuaderno de primeras palabras: casi todas las actividades
% terminan coloreando o dibujando, y el sitio es todo el que haya. Lo que
% tiene que ir abajo del todo (la línea de escribir, el cartel de cada
% diez páginas) va en la parte de abajo de la caja (\tcblower, estilo
% `abajo`); lo que tiene que ir centrado en el hueco, con `centrado abajo`.
\tcbset{
  cajabase/.style={
    enhanced,
    boxrule=1.1pt,
    arc=3mm,
    left=6mm, right=6mm, top=5mm, bottom=5mm,
    fonttitle=\bfseries\sffamily\large,
    coltitle=white,
  },
  cajadia/.style={cajabase, height fill},
  abajo/.style={space to upper, segmentation hidden, middle=2mm},
  centrado abajo/.style={space to lower, valign lower=center, segmentation hidden, middle=2mm},
}

\newtcolorbox{cajaNumero}{cajabase, top=3mm, bottom=3mm,
  colback=colorLecturaClaro, colframe=colorLectura, colbacktitle=colorLectura,
  title=\lblNumeroDeHoy}

\newtcolorbox{cajaTraza}[1][]{cajadia, #1,
  colback=colorRespondeClaro, colframe=colorResponde, colbacktitle=colorResponde,
  title=\lblTraza}
\newtcolorbox{cajaColorea}[1][]{cajadia, #1,
  colback=colorCompletaClaro, colframe=colorCompleta, colbacktitle=colorCompleta,
  title=\lblColorea}
\newtcolorbox{cajaRodea}[1][]{cajadia, #1,
  colback=colorRelacionaClaro, colframe=colorRelaciona, colbacktitle=colorRelaciona,
  title=\lblRodea}
\newtcolorbox{cajaCuenta}[1][]{cajadia, #1,
  colback=colorRelacionaClaro, colframe=colorRelaciona, colbacktitle=colorRelaciona,
  title=\lblCuenta}
\newtcolorbox{cajaBusca}[1][]{cajadia, #1,
  colback=colorCompletaClaro, colframe=colorCompleta, colbacktitle=colorCompleta,
  title=\lblBusca}
\newtcolorbox{cajaUne}[1][]{cajadia, #1,
  colback=colorRelacionaClaro, colframe=colorRelaciona, colbacktitle=colorRelaciona,
  title=\lblUne}
\newtcolorbox{cajaCompleta}[1][]{cajadia, #1,
  colback=colorRespondeClaro, colframe=colorResponde, colbacktitle=colorResponde,
  title=\lblCompleta}
% Las del invierno.
\newtcolorbox{cajaDecena}[1][]{cajadia, #1,
  colback=colorRelacionaClaro, colframe=colorRelaciona, colbacktitle=colorRelaciona,
  title=\lblDecena}
\newtcolorbox{cajaCompara}[1][]{cajadia, #1,
  colback=colorRelacionaClaro, colframe=colorRelaciona, colbacktitle=colorRelaciona,
  title=\lblCompara}
\newtcolorbox{cajaVecinos}[1][]{cajadia, #1,
  colback=colorRespondeClaro, colframe=colorResponde, colbacktitle=colorResponde,
  title=\lblVecinos}
\newtcolorbox{cajaRecta}[1][]{cajadia, #1,
  colback=colorRespondeClaro, colframe=colorResponde, colbacktitle=colorResponde,
  title=\lblRecta}
\newtcolorbox{cajaOrdena}[1][]{cajadia, #1,
  colback=colorCompletaClaro, colframe=colorCompleta, colbacktitle=colorCompleta,
  title=\lblOrdena}
\newtcolorbox{cajaJunta}[1][]{cajadia, #1,
  colback=colorDibujaClaro, colframe=colorDibuja, colbacktitle=colorDibuja,
  title=\lblJunta}
\newtcolorbox{cajaDibuja}[1][]{cajadia, #1,
  colback=colorDibujaClaro, colframe=colorDibuja, colbacktitle=colorDibuja,
  title=\lblDibuja}
\newtcolorbox{cajaRepasa}[1][]{cajadia, #1,
  colback=colorCreaClaro, colframe=colorCrea, colbacktitle=colorCrea,
  title=\lblRepasa}

% La instrucción de cada actividad: la lee el adulto.
\newcommand{\instruccion}[1]{{\footnotesize\color{colorGris}#1\par}}
% El enunciado grande, lo que se hace (lo que el niño o la niña puede
% seguir con el dedo): "Colorea 2 pelotas."
\newcommand{\enunciado}[1]{{\LARGE #1\par}}
% Un renglón para escribir el número: el número de muestra, en gris, y
% sitio en la misma línea para escribirlo tres o cuatro veces más.
\newcommand{\renglonTraza}[1]{%
  \par\noindent\begin{tikzpicture}
    \draw[line width=0.6pt] (0,0) -- (\linewidth-1pt,0);
    \node[anchor=base west, inner sep=0pt, text=colorGris!55,
      font=\fontsize{44}{44}\selectfont\bfseries] at (2mm,0) {#1};
  \end{tikzpicture}\par\vspace{12mm}}
% El hueco en el que se escribe la respuesta: "10 y 3 son [ ]".
\newcommand{\huecoRespuesta}{\raisebox{-5mm}{\tikz\draw[line width=1.1pt,
  rounded corners=2mm, fill=white] (0,0) rectangle (26mm,18mm);}}
% Lo que no cabe en su sitio, encogido hasta que quepa (#1, el ancho):
% los nombres largos, como "diecisiete", en la caja de arriba.
\newsavebox{\cajaAjustada}
\newcommand{\ajustado}[2]{\sbox{\cajaAjustada}{#2}%
  \ifdim\wd\cajaAjustada>#1\relax\resizebox{#1}{!}{\usebox{\cajaAjustada}}%
  \else\usebox{\cajaAjustada}\fi}
% El cartel de cada diez páginas, al pie de "Repasa".
\newcommand{\cartel}[1]{\begin{center}\bfseries\Large\color{colorCrea}#1\end{center}}

\newcommand{\casilla}{\raisebox{-0.15ex}{\framebox[3.4mm][c]{\rule{0pt}{2.6mm}}}}
\newenvironment{listaRepaso}{%
  \begin{itemize}[label=\casilla, leftmargin=8mm, itemsep=2mm, topsep=1mm]\large
}{%
  \end{itemize}
}

% =========================================================================
% La página de un día
% =========================================================================
% La cabecera es la de "Aprendo a leer": el día, la semana y el
% trimestre; el tema de la semana (el mismo que esa semana en "Aprendo a
% leer", para que quien lleve los dos cuadernos cuente lo que lee), y los
% campos para el adulto -- aquí "Hecho" en vez de "Leído". Deja
% \label{dia:N}: tools/check_pages.py comprueba con esas etiquetas que
% cada día ocupa exactamente una página.
\newcommand{\cabeceraDia}[4]{%
  \label{dia:#1}%
  \noindent{\LARGE\bfseries\lblDia~#1}\quad
  {\large(\lblSemana~#2 \textperiodcentered\ \lblTrimestre~#3)}\par
  {\normalsize\bfseries\color{colorGris}#4}\par
  \vspace{1mm}
  \textbf{\lblFecha:} \rule{22mm}{0.4pt}\hspace{4mm}%
  \textbf{\lblHecho:} \casilla\ \lblHechoSola\quad
  \casilla\ \lblHechoConAyuda\quad
  \casilla\ \lblHechoConDificultad\hspace{4mm}%
  \textbf{\lblNotas:} \rule{20mm}{0.4pt}\par
  \vspace{3mm}
}

\makeatletter
\def\ps@pienumeros{%
  \let\@oddhead\@empty\let\@evenhead\@empty
  \def\@oddfoot{\hfil\footnotesize\color{colorGris}\lblDia~\diapagina@actual~/ \totaldias\hfil}%
  \let\@evenfoot\@oddfoot
}
\newenvironment{diapagina}[4]{%
  \gdef\diapagina@actual{#1}%
  \thispagestyle{pienumeros}%
  \cabeceraDia{#1}{#2}{#3}{#4}%
}{%
  \par\newpage
}

% Una etiqueta que solo guarda la página en la que cae, para
% tools/check_pages.py: la clave de respuestas empieza con
% \etiquetaPagina{clave}, así se comprueba también que el último día no
% se ha desbordado a una segunda página.
\newcommand{\etiquetaPagina}[1]{\begingroup\def\@currentlabel{}\label{#1}\endgroup}
\makeatother

% =========================================================================
% El número de hoy
% =========================================================================
% El marco de diez: dos filas de cinco casillas, con un punto en las N
% primeras -- la manera de ver de un vistazo cuánto es un número hasta
% el 10 (5 es una fila entera; 7, una fila y dos más).
\newcommand{\marcoDiez}[1]{%
  \begin{tikzpicture}[x=7.5mm, y=7.5mm, line width=0.9pt]
    \foreach \i in {0,...,9} {
      \pgfmathtruncatemacro\col{mod(\i,5)}
      \pgfmathtruncatemacro\fila{\i<5 ? 1 : 0}
      \draw (\col,\fila) rectangle ++(1,1);
      \ifnum\i<#1\relax
        \fill[colorLectura] (\col+0.5,\fila+0.5) circle (0.32);
      \fi
    }
  \end{tikzpicture}}
% Del 11 al 20, dos marcos, uno encima de otro: el primero, lleno.
\newcommand{\marcosDiez}[1]{%
  \ifnum#1>10\relax
    \marcoDiez{10}\par\vspace{1.5mm}\marcoDiez{\numexpr#1-10\relax}%
  \else
    \marcoDiez{#1}%
  \fi}

% La caja de arriba: el número grande, su nombre, el marco de diez y las
% cosas que se cuentan ese día (#3: un tikzpicture que compone
% tools/gen_numeros.py), y debajo la frase de la historia, que lee el
% adulto. La columna del número mide 28 mm, o lo que mida el número si
% tiene dos cifras (el 10): no se encoge, y lo que se estrecha es la de
% las cosas (tools/gen_numeros.py las dibuja más pequeñas).
\newlength{\anchoNumeroDeHoy}
\newcommand{\fuenteNumeroDeHoy}{\fontsize{92}{92}\selectfont\bfseries}
\newcommand{\numeroDeHoy}[4]{%
  \settowidth{\anchoNumeroDeHoy}{\fuenteNumeroDeHoy #1}%
  \ifdim\anchoNumeroDeHoy<28mm \setlength{\anchoNumeroDeHoy}{28mm}\fi
  \begin{cajaNumero}
  \begin{tabularx}{\linewidth}{@{}>{\centering\arraybackslash}m{\anchoNumeroDeHoy}@{\hspace{3mm}}>{\centering\arraybackslash}m{42mm}@{\hspace{3mm}}>{\centering\arraybackslash}X@{}}
  {\fuenteNumeroDeHoy\color{colorLectura}#1} &
  \ajustado{42mm}{\fontsize{30}{34}\selectfont\bfseries #2}\par\vspace{3mm}\marcosDiez{#1} &
  #3 \\
  \end{tabularx}
  \par\vspace{2mm}
  {\large #4\par}
  \end{cajaNumero}\par\vspace{4mm}}

% =========================================================================
% Medalla de fin de trimestre (tools/gen_numeros.py, PLANTILLA_MEDALLA)
% =========================================================================
\newcommand{\medalla}[1]{%
  \begin{tikzpicture}[line width=1.6pt, color=colorCrea]
    \draw (0,0) circle (1.1);
    \draw (0,0) circle (0.85);
    \node at (0,0) {\Large\bfseries #1};
    \draw (-0.35,-1.05) -- (-0.7,-2.2) -- (-0.15,-1.75) -- cycle;
    \draw (0.35,-1.05)  -- (0.7,-2.2)  -- (0.15,-1.75)  -- cycle;
  \end{tikzpicture}}

% =========================================================================
% La clave de respuestas
% =========================================================================
\newcommand{\claveEntrada}[3]{%
  \item \textbf{\lblDia\ #1} \textcolor{colorGris}{(#2)}: #3%
}
, y los mismos nombres de color
% apuntan a negro y gris: la misma página, sin tinta de color.
\makeatletter
\newif\ifbwbook
\def\aln@bookcolor@bw{bw}
\@ifundefined{bookcolor}{\bwbookfalse}{%
  \ifx\bookcolor\aln@bookcolor@bw \bwbooktrue \else \bwbookfalse \fi
}
\makeatother

\ifbwbook
  \definecolor{colorLectura}{gray}{0}
  \definecolor{colorLecturaClaro}{gray}{0.95}
  \definecolor{colorDibuja}{gray}{0}
  \definecolor{colorDibujaClaro}{gray}{0.95}
  \definecolor{colorCompleta}{gray}{0}
  \definecolor{colorCompletaClaro}{gray}{0.95}
  \definecolor{colorResponde}{gray}{0}
  \definecolor{colorRespondeClaro}{gray}{0.95}
  \definecolor{colorRelaciona}{gray}{0}
  \definecolor{colorRelacionaClaro}{gray}{0.95}
  \definecolor{colorCrea}{gray}{0}
  \definecolor{colorCreaClaro}{gray}{0.95}
  \definecolor{colorGris}{gray}{0.38}
  \definecolor{colorPunto}{gray}{0.25}
\else
  \definecolor{colorLectura}{HTML}{2E7D32}
  \definecolor{colorLecturaClaro}{HTML}{EAF6EC}
  \definecolor{colorDibuja}{HTML}{1565C0}
  \definecolor{colorDibujaClaro}{HTML}{E7F0FB}
  \definecolor{colorCompleta}{HTML}{C2185B}
  \definecolor{colorCompletaClaro}{HTML}{FBE9EF}
  \definecolor{colorResponde}{HTML}{6A1B9A}
  \definecolor{colorRespondeClaro}{HTML}{F2E9F7}
  \definecolor{colorRelaciona}{HTML}{E65100}
  \definecolor{colorRelacionaClaro}{HTML}{FDEEE2}
  \definecolor{colorCrea}{HTML}{AD1457}
  \definecolor{colorCreaClaro}{HTML}{FBE7EE}
  \definecolor{colorGris}{HTML}{616161}
  \definecolor{colorPunto}{HTML}{6A1B9A}
\fi

% =========================================================================
% Los dibujos
% =========================================================================
% diagrams/kit.tex: las piezas de los dibujos de «First Words» (de
% "Aprendo a leer"), con las que están hechas muchas de las cosas que se
% cuentan; diagrams/objetos.tex: esas cosas, cada una en la misma caja.
\input{diagrams/kit}
% diagrams/objetos.tex -- las cosas que se cuentan: cada una dibujada
% dentro de la misma caja, de -1 a 1 en x y en y, para que
% tools/gen_numeros.py las pueda poner en fila, en grupos o en el marco
% de diez sin saber qué es cada una (ver OBJETOS en tools/gen_numeros.py,
% que dice qué macro dibuja cada cosa y cómo se llama en singular y en
% plural). Todas las zonas quedan cerradas y rellenas de blanco, para
% que cada una se pueda colorear.
%
% Muchas son los dibujos de «First Words» (diagrams/kit.tex, del
% repositorio de "Aprendo a leer"), encogidos o movidos hasta caber en la
% caja; las demás son nuevas y siguen el mismo estilo de línea.
\tikzset{objeto/.style={line width=1.1pt, line cap=round, line join=round}}

% Una manzana, con el rabito y una hoja.
\newcommand{\objManzana}{%
  \filldraw[fill=white] (0,0.55) .. controls (-0.35,0.9) and (-1.0,0.78) .. (-0.95,0.02)
    .. controls (-0.9,-0.72) and (-0.35,-1.0) .. (0,-0.8)
    .. controls (0.35,-1.0) and (0.9,-0.72) .. (0.95,0.02)
    .. controls (1.0,0.78) and (0.35,0.9) .. (0,0.55) -- cycle;
  \draw (0,0.55) .. controls (0.02,0.8) .. (0.1,0.98);
  \filldraw[fill=white] (0.08,0.85) .. controls (0.32,1.05) and (0.62,0.98) .. (0.7,0.84)
    .. controls (0.5,0.68) and (0.25,0.7) .. (0.08,0.85) -- cycle;}

% Una pelota de playa, con cuatro gajos.
\newcommand{\objPelota}{%
  \filldraw[fill=white] (0,0) circle (0.9);
  \draw (-0.9,0) .. controls (-0.3,0.38) and (0.3,0.38) .. (0.9,0);
  \draw (0,0.9) .. controls (-0.38,0.3) and (-0.38,-0.3) .. (0,-0.9);}

% Un hueso, tumbado: la barra y los cuatro nudos, un solo contorno.
\newcommand{\objHueso}{%
  \filldraw[fill=white] (-0.492,0.16) -- (0.492,0.16)
    arc[start angle=197.2, end angle=-62.7, radius=0.27]
    arc[start angle=62.7, end angle=-197.2, radius=0.27]
    -- (-0.492,-0.16)
    arc[start angle=17.2, end angle=-242.7, radius=0.27]
    arc[start angle=242.7, end angle=-17.2, radius=0.27] -- cycle;}

% Un sol: el círculo y ocho rayos.
\newcommand{\objSol}{%
  \filldraw[fill=white] (0,0) circle (0.5);
  \foreach \a in {0,45,...,315} {\draw (\a:0.66) -- (\a:0.96);}}

% Una huella de Toby.
\newcommand{\objHuella}{%
  \begin{scope}[shift={(0,-0.28)}, scale=1.65] \huellaFW \end{scope}}

% Toby, sentado.
\newcommand{\objToby}{%
  \begin{scope}[shift={(0.38,-0.88)}, scale=0.4] \tobySentadoFW \end{scope}}

% Un globo, con su hilo.
\newcommand{\objGlobo}{%
  \begin{scope}[shift={(0,-0.85)}, scale=0.72] \globoFW{0.2}{-0.2} \end{scope}}

% Una estrella.
\newcommand{\objEstrella}{%
  \begin{scope}[scale=0.95] \estrellaFW \end{scope}}

% Una hoja de otoño.
\newcommand{\objHoja}{%
  \begin{scope}[shift={(0,-0.8)}, scale=1.2] \hojaFW \end{scope}}

% Un corazón.
\newcommand{\objCorazon}{%
  \begin{scope}[shift={(0,-0.82)}, scale=1.35] \corazonFW \end{scope}}

% Un caramelo.
\newcommand{\objCaramelo}{%
  \begin{scope}[scale=1.25] \carameloFW \end{scope}}

% Un pez.
\newcommand{\objPez}{%
  \begin{scope}[shift={(0.2,0)}, scale=0.72] \pezFW \end{scope}}


% =========================================================================
% Cajas
% =========================================================================
% La caja de actividad llena todo lo que queda de página (`height fill`),
% como en el cuaderno de primeras palabras: casi todas las actividades
% terminan coloreando o dibujando, y el sitio es todo el que haya. Lo que
% tiene que ir abajo del todo (la línea de escribir, el cartel de cada
% diez páginas) va en la parte de abajo de la caja (\tcblower, estilo
% `abajo`); lo que tiene que ir centrado en el hueco, con `centrado abajo`.
\tcbset{
  cajabase/.style={
    enhanced,
    boxrule=1.1pt,
    arc=3mm,
    left=6mm, right=6mm, top=5mm, bottom=5mm,
    fonttitle=\bfseries\sffamily\large,
    coltitle=white,
  },
  cajadia/.style={cajabase, height fill},
  abajo/.style={space to upper, segmentation hidden, middle=2mm},
  centrado abajo/.style={space to lower, valign lower=center, segmentation hidden, middle=2mm},
}

\newtcolorbox{cajaNumero}{cajabase, top=3mm, bottom=3mm,
  colback=colorLecturaClaro, colframe=colorLectura, colbacktitle=colorLectura,
  title=\lblNumeroDeHoy}

\newtcolorbox{cajaTraza}[1][]{cajadia, #1,
  colback=colorRespondeClaro, colframe=colorResponde, colbacktitle=colorResponde,
  title=\lblTraza}
\newtcolorbox{cajaColorea}[1][]{cajadia, #1,
  colback=colorCompletaClaro, colframe=colorCompleta, colbacktitle=colorCompleta,
  title=\lblColorea}
\newtcolorbox{cajaRodea}[1][]{cajadia, #1,
  colback=colorRelacionaClaro, colframe=colorRelaciona, colbacktitle=colorRelaciona,
  title=\lblRodea}
\newtcolorbox{cajaCuenta}[1][]{cajadia, #1,
  colback=colorRelacionaClaro, colframe=colorRelaciona, colbacktitle=colorRelaciona,
  title=\lblCuenta}
\newtcolorbox{cajaBusca}[1][]{cajadia, #1,
  colback=colorCompletaClaro, colframe=colorCompleta, colbacktitle=colorCompleta,
  title=\lblBusca}
\newtcolorbox{cajaUne}[1][]{cajadia, #1,
  colback=colorRelacionaClaro, colframe=colorRelaciona, colbacktitle=colorRelaciona,
  title=\lblUne}
\newtcolorbox{cajaCompleta}[1][]{cajadia, #1,
  colback=colorRespondeClaro, colframe=colorResponde, colbacktitle=colorResponde,
  title=\lblCompleta}
% Las del invierno.
\newtcolorbox{cajaDecena}[1][]{cajadia, #1,
  colback=colorRelacionaClaro, colframe=colorRelaciona, colbacktitle=colorRelaciona,
  title=\lblDecena}
\newtcolorbox{cajaCompara}[1][]{cajadia, #1,
  colback=colorRelacionaClaro, colframe=colorRelaciona, colbacktitle=colorRelaciona,
  title=\lblCompara}
\newtcolorbox{cajaVecinos}[1][]{cajadia, #1,
  colback=colorRespondeClaro, colframe=colorResponde, colbacktitle=colorResponde,
  title=\lblVecinos}
\newtcolorbox{cajaRecta}[1][]{cajadia, #1,
  colback=colorRespondeClaro, colframe=colorResponde, colbacktitle=colorResponde,
  title=\lblRecta}
\newtcolorbox{cajaOrdena}[1][]{cajadia, #1,
  colback=colorCompletaClaro, colframe=colorCompleta, colbacktitle=colorCompleta,
  title=\lblOrdena}
\newtcolorbox{cajaJunta}[1][]{cajadia, #1,
  colback=colorDibujaClaro, colframe=colorDibuja, colbacktitle=colorDibuja,
  title=\lblJunta}
\newtcolorbox{cajaDibuja}[1][]{cajadia, #1,
  colback=colorDibujaClaro, colframe=colorDibuja, colbacktitle=colorDibuja,
  title=\lblDibuja}
\newtcolorbox{cajaRepasa}[1][]{cajadia, #1,
  colback=colorCreaClaro, colframe=colorCrea, colbacktitle=colorCrea,
  title=\lblRepasa}

% La instrucción de cada actividad: la lee el adulto.
\newcommand{\instruccion}[1]{{\footnotesize\color{colorGris}#1\par}}
% El enunciado grande, lo que se hace (lo que el niño o la niña puede
% seguir con el dedo): "Colorea 2 pelotas."
\newcommand{\enunciado}[1]{{\LARGE #1\par}}
% Un renglón para escribir el número: el número de muestra, en gris, y
% sitio en la misma línea para escribirlo tres o cuatro veces más.
\newcommand{\renglonTraza}[1]{%
  \par\noindent\begin{tikzpicture}
    \draw[line width=0.6pt] (0,0) -- (\linewidth-1pt,0);
    \node[anchor=base west, inner sep=0pt, text=colorGris!55,
      font=\fontsize{44}{44}\selectfont\bfseries] at (2mm,0) {#1};
  \end{tikzpicture}\par\vspace{12mm}}
% El hueco en el que se escribe la respuesta: "10 y 3 son [ ]".
\newcommand{\huecoRespuesta}{\raisebox{-5mm}{\tikz\draw[line width=1.1pt,
  rounded corners=2mm, fill=white] (0,0) rectangle (26mm,18mm);}}
% Lo que no cabe en su sitio, encogido hasta que quepa (#1, el ancho):
% los nombres largos, como "diecisiete", en la caja de arriba.
\newsavebox{\cajaAjustada}
\newcommand{\ajustado}[2]{\sbox{\cajaAjustada}{#2}%
  \ifdim\wd\cajaAjustada>#1\relax\resizebox{#1}{!}{\usebox{\cajaAjustada}}%
  \else\usebox{\cajaAjustada}\fi}
% El cartel de cada diez páginas, al pie de "Repasa".
\newcommand{\cartel}[1]{\begin{center}\bfseries\Large\color{colorCrea}#1\end{center}}

\newcommand{\casilla}{\raisebox{-0.15ex}{\framebox[3.4mm][c]{\rule{0pt}{2.6mm}}}}
\newenvironment{listaRepaso}{%
  \begin{itemize}[label=\casilla, leftmargin=8mm, itemsep=2mm, topsep=1mm]\large
}{%
  \end{itemize}
}

% =========================================================================
% La página de un día
% =========================================================================
% La cabecera es la de "Aprendo a leer": el día, la semana y el
% trimestre; el tema de la semana (el mismo que esa semana en "Aprendo a
% leer", para que quien lleve los dos cuadernos cuente lo que lee), y los
% campos para el adulto -- aquí "Hecho" en vez de "Leído". Deja
% \label{dia:N}: tools/check_pages.py comprueba con esas etiquetas que
% cada día ocupa exactamente una página.
\newcommand{\cabeceraDia}[4]{%
  \label{dia:#1}%
  \noindent{\LARGE\bfseries\lblDia~#1}\quad
  {\large(\lblSemana~#2 \textperiodcentered\ \lblTrimestre~#3)}\par
  {\normalsize\bfseries\color{colorGris}#4}\par
  \vspace{1mm}
  \textbf{\lblFecha:} \rule{22mm}{0.4pt}\hspace{4mm}%
  \textbf{\lblHecho:} \casilla\ \lblHechoSola\quad
  \casilla\ \lblHechoConAyuda\quad
  \casilla\ \lblHechoConDificultad\hspace{4mm}%
  \textbf{\lblNotas:} \rule{20mm}{0.4pt}\par
  \vspace{3mm}
}

\makeatletter
\def\ps@pienumeros{%
  \let\@oddhead\@empty\let\@evenhead\@empty
  \def\@oddfoot{\hfil\footnotesize\color{colorGris}\lblDia~\diapagina@actual~/ \totaldias\hfil}%
  \let\@evenfoot\@oddfoot
}
\newenvironment{diapagina}[4]{%
  \gdef\diapagina@actual{#1}%
  \thispagestyle{pienumeros}%
  \cabeceraDia{#1}{#2}{#3}{#4}%
}{%
  \par\newpage
}

% Una etiqueta que solo guarda la página en la que cae, para
% tools/check_pages.py: la clave de respuestas empieza con
% \etiquetaPagina{clave}, así se comprueba también que el último día no
% se ha desbordado a una segunda página.
\newcommand{\etiquetaPagina}[1]{\begingroup\def\@currentlabel{}\label{#1}\endgroup}
\makeatother

% =========================================================================
% El número de hoy
% =========================================================================
% El marco de diez: dos filas de cinco casillas, con un punto en las N
% primeras -- la manera de ver de un vistazo cuánto es un número hasta
% el 10 (5 es una fila entera; 7, una fila y dos más).
\newcommand{\marcoDiez}[1]{%
  \begin{tikzpicture}[x=7.5mm, y=7.5mm, line width=0.9pt]
    \foreach \i in {0,...,9} {
      \pgfmathtruncatemacro\col{mod(\i,5)}
      \pgfmathtruncatemacro\fila{\i<5 ? 1 : 0}
      \draw (\col,\fila) rectangle ++(1,1);
      \ifnum\i<#1\relax
        \fill[colorLectura] (\col+0.5,\fila+0.5) circle (0.32);
      \fi
    }
  \end{tikzpicture}}
% Del 11 al 20, dos marcos, uno encima de otro: el primero, lleno.
\newcommand{\marcosDiez}[1]{%
  \ifnum#1>10\relax
    \marcoDiez{10}\par\vspace{1.5mm}\marcoDiez{\numexpr#1-10\relax}%
  \else
    \marcoDiez{#1}%
  \fi}

% La caja de arriba: el número grande, su nombre, el marco de diez y las
% cosas que se cuentan ese día (#3: un tikzpicture que compone
% tools/gen_numeros.py), y debajo la frase de la historia, que lee el
% adulto. La columna del número mide 28 mm, o lo que mida el número si
% tiene dos cifras (el 10): no se encoge, y lo que se estrecha es la de
% las cosas (tools/gen_numeros.py las dibuja más pequeñas).
\newlength{\anchoNumeroDeHoy}
\newcommand{\fuenteNumeroDeHoy}{\fontsize{92}{92}\selectfont\bfseries}
\newcommand{\numeroDeHoy}[4]{%
  \settowidth{\anchoNumeroDeHoy}{\fuenteNumeroDeHoy #1}%
  \ifdim\anchoNumeroDeHoy<28mm \setlength{\anchoNumeroDeHoy}{28mm}\fi
  \begin{cajaNumero}
  \begin{tabularx}{\linewidth}{@{}>{\centering\arraybackslash}m{\anchoNumeroDeHoy}@{\hspace{3mm}}>{\centering\arraybackslash}m{42mm}@{\hspace{3mm}}>{\centering\arraybackslash}X@{}}
  {\fuenteNumeroDeHoy\color{colorLectura}#1} &
  \ajustado{42mm}{\fontsize{30}{34}\selectfont\bfseries #2}\par\vspace{3mm}\marcosDiez{#1} &
  #3 \\
  \end{tabularx}
  \par\vspace{2mm}
  {\large #4\par}
  \end{cajaNumero}\par\vspace{4mm}}

% =========================================================================
% Medalla de fin de trimestre (tools/gen_numeros.py, PLANTILLA_MEDALLA)
% =========================================================================
\newcommand{\medalla}[1]{%
  \begin{tikzpicture}[line width=1.6pt, color=colorCrea]
    \draw (0,0) circle (1.1);
    \draw (0,0) circle (0.85);
    \node at (0,0) {\Large\bfseries #1};
    \draw (-0.35,-1.05) -- (-0.7,-2.2) -- (-0.15,-1.75) -- cycle;
    \draw (0.35,-1.05)  -- (0.7,-2.2)  -- (0.15,-1.75)  -- cycle;
  \end{tikzpicture}}

% =========================================================================
% La clave de respuestas
% =========================================================================
\newcommand{\claveEntrada}[3]{%
  \item \textbf{\lblDia\ #1} \textcolor{colorGris}{(#2)}: #3%
}
